\subsection{总体分析}
四大家鱼回升、江豚回暖和长江鲟产卵场重现都指向禁渔后的恢复；“秃海”、通道受阻和入侵扩散又说明，单项资源量不足以解释收益在不同营养级和不同外部压力下的去向。我们因此不把注意力停在单项资源量上，分析转向收益所处的层级和压力来源。

五问彼此关联。问题一定位恢复发生在哪一层，问题二考察资源端的变化能否传到珍稀物种；问题三把短期响应压缩成Hastings-Powell食物链，问题四判断复杂轨道会不会随初值变化；问题五把污染和入侵加入已有状态方程，比较剩余收益。全文因此把“恢复是否发生”与“恢复是否稳得住、是否算健康”分别落在对应的计算环节。

